\chapter{总结与展望}

\section{研究工作总结}

交通资产清查是支撑道路基础设施数字化台账更新、风险研判与精细化养护的重要基础环节。针对现有方法在复杂道路场景下面临的“资产发现不完整”与“异常状态记录不充分”两类核心问题，本文围绕交通资产清查任务，从主动感知检测与知识引导属性识别增强两个层面展开研究，构建了面向交通资产清查的完整方法链路。全文的研究目标在于：在保证工程可实现性的前提下，提高交通资产目标发现的完整性、异常状态识别的准确性以及结果输出的结构化程度。

在数据与任务建模层面，本文首先基于 WHU-Infra3D 数据集建立统一实验基础。针对全景影像几何畸变对检测任务的不利影响，本文将上海与武汉子集的等距柱状投影全景图进行四分透视投影切割，由单幅 ERP 全景图生成 4 幅透视图像，并据此构建跨城市检测实验数据。其中，武汉训练集共包含 6712 张四分投影视图，上海测试集包含 4700 张四分投影视图。对于属性识别任务，本文进一步采用“距离曝光点 20\,m 筛选 + 标注框 1.5 倍外扩裁剪”的方式构建 7311 张真实局部样本，并在此基础上生成 4197 条异常属性样本，形成统一的训练与测试数据组织方案，为后续检测与识别实验提供了统一的数据支撑。

在交通资产检测方面，本文提出了基于主动感知的多智能体三阶段检测框架。该框架通过 Detection Agent 与 Crop Agent 的协同，将传统单次全图检测重构为“全局初检—候选区域推荐—局部精检—结果融合”的主动感知流程。与穷举式规则裁剪不同，本文方法强调在有限局部检测预算下主动筛选更值得复查的区域，从而提升对潜在漏检目标的恢复能力。实验结果表明，所提框架对不同检测器具有较强适配性，在 RexOmni、DINO、Faster R-CNN、DEIM、Grounding DINO 和 YOLO26 六类异构检测器上均取得稳定增益，Macro-F1 与 Micro-F1 平均分别提升 8.18\% 和 6.41\%。进一步分析发现，本文方法的主要收益集中体现在小目标、密集区域与复杂局部场景中，说明主动感知机制能够有效缓解全图检测中普遍存在的尺度压缩与局部遗漏问题。

在交通资产属性识别方面，本文提出了领域知识引导的长尾数据生成与增强方法。该方法将交通行业规范、运维经验与视觉先验转化为可控 Prompt 约束，面向锈蚀、风化老旧、结构损坏、井盖移位及垃圾满溢等关键异常状态开展定向补样，并在统一多模态识别框架下验证其对结构化属性识别的促进作用。实验结果表明，长尾增强模型在统一测试集上的样本级严格准确率、属性级准确率和异常属性准确率分别达到 88.78\%、96.12\% 和 93.72\%，相较仅使用真实样本训练的模型分别提升 30.82、11.60 和 92.71 个百分点，正常--异常性能差距由 95.63\% 显著缩小至 6.08\%。这表明，知识引导长尾增强不仅改善了模型对关键低频异常字段的识别能力，也提升了多字段联合输出的一致性与稳定性。

总体而言，本文围绕“查全资产”与“看懂状态”两个关键环节，形成了从数据构建、任务建模、方法设计到实验验证的较完整研究闭环。一方面，主动感知检测方法提升了交通资产发现的完整性；另一方面，知识引导长尾增强方法提升了异常属性识别的可靠性。二者共同构成了面向交通资产清查任务的技术支撑体系，也为后续开展更高精度、更强泛化能力的交通资产智能巡检研究奠定了基础。

\section{创新点归纳}

在前述研究工作的基础上，本文的创新点可进一步归纳为以下三个方面：

\begin{enumerate}
    \item \textbf{面向交通资产漏检问题，提出了主动感知驱动的三阶段检测框架。} 针对传统全图检测在小目标、密集目标和复杂局部场景中易发生漏检的局限，本文将检测流程重构为“全局初检—局部复查—结果融合”的主动感知范式，使检测过程由被动单次识别转变为面向高价值区域的任务驱动式复查，从方法机制上增强了对潜在漏检目标的恢复能力。
    \item \textbf{面向交通资产异常属性长尾问题，提出了领域知识引导的样本构建与增强方法。} 针对关键异常样本稀缺、真实采集成本高且分布不均衡的问题，本文将行业规范、运维经验与视觉先验前移为训练分布塑造机制，构建了面向异常状态的知识约束生成方法，实现了对锈蚀、结构损坏、井盖移位和垃圾满溢等低频异常属性的定向补样，为结构化属性识别提供了更具针对性的训练支撑。
    \item \textbf{围绕交通资产清查任务，构建了统一的数据组织与实验验证链路。} 本文在统一数据基础、统一任务组织形式与统一评测口径下，系统验证了主动感知检测与知识引导长尾增强两项方法在交通资产清查场景中的有效性与工程潜力，从而提升了研究结论的可比性、完整性与应用说服力。
\end{enumerate}

上述创新点并非彼此孤立，而是分别对应交通资产清查中的“目标发现”与“状态理解”两个关键阶段，并在统一实验体系下形成相互支撑的整体技术方案。这也是本文相较于单一检测研究或单一识别研究的一个重要特点。

\section{展望}

尽管本文围绕交通资产清查任务开展了系统研究，并在检测与属性识别两个层面取得了一定进展，但仍存在若干值得进一步完善和深入研究的问题。未来的研究工作可从以下几个方面继续展开。

首先，在主动感知检测方面，当前三阶段框架虽然能够稳定提升整体检测性能，但其候选区域推荐质量仍然依赖第一阶段初检所提供的先验线索。当原始检测器对某些极端困难目标完全缺乏感知，或场景中存在严重遮挡、极端模糊和异常复杂背景时，Crop Agent 对高价值区域的筛选能力仍可能受到限制。此外，本文目前主要围绕二维透视图像开展检测实验，尚未充分利用 WHU-Infra3D 数据集中点云、深度与多模态几何信息，这使得方法在三维空间一致性建模和跨模态互补方面仍有提升空间。未来可进一步探索基于不确定性估计的候选区域预算分配、自适应多轮主动感知策略，以及结合点云几何约束的跨模态资产检测方法。

其次，在属性识别与长尾增强方面，本文所构建的知识引导生成方法虽然显著改善了关键异常字段的识别效果，但生成样本与真实异常场景之间仍可能存在分布偏差。尤其对于结构性破坏、复合异常和强场景耦合异常，当前基于局部编辑的生成方式仍难以完全复现真实道路环境中的复杂视觉表现。此外，统一测试集中的一部分异常样本仍来源于生成数据，因此对模型真实泛化能力的评估仍有进一步加强的必要。未来可从两个方向进一步推进：一方面，持续补充真实异常回采样本，构建覆盖更广、真实性更强的异常属性数据集；另一方面，将领域知识约束、自动质量筛选与人工复核机制进一步结合，提升生成数据的保真度与工程可用性。

再次，从系统集成角度看，本文目前仍将检测与属性识别视为前后串联的两个阶段，尚未充分研究两者之间更紧密的协同关系。例如，检测阶段发现的候选区域置信度、空间上下文与多目标关系，可能为属性识别提供额外先验；反过来，属性识别得到的异常状态信息也可能用于反向修正检测结果、辅助资产优先级排序与风险处置决策。未来可进一步探索面向交通资产清查的一体化建模框架，将资产发现、属性识别、异常研判与结构化记录纳入统一任务链条，以提升系统级性能与应用闭环能力。

最后，在工程部署与应用推广层面，本文的实验验证主要基于跨城市离线评测场景展开，对于实时处理约束、设备算力限制、长周期数据更新机制以及人工复核流程耦合等实际问题仍缺乏系统研究。未来可在以下几个方面进一步深化：其一，面向车载或路侧设备开展轻量化部署与推理效率优化；其二，构建“模型识别—人工复核—数据回流—持续训练”的闭环更新机制；其三，将本文方法与交通资产管理平台、巡检台账系统和养护决策流程进一步衔接，验证其在真实业务链路中的稳定性与长期价值。

综上，本文的研究工作为交通资产清查中的主动感知检测与知识引导属性识别提供了初步而系统的技术基础，但在多模态融合、真实异常数据积累、端到端协同建模以及工程闭环应用等方面仍存在广阔的拓展空间。后续研究若能在上述方向持续推进，将有望进一步提升交通资产智能清查系统的准确性、鲁棒性与实际落地能力。
